% Part: first-order-logic
% Chapter: beyond
% Section: second-order-logic

\documentclass[../../../include/open-logic-section]{subfiles}

\begin{document}

\olfileid[pt]{fol}{byd}{sol}

\olsection{Lógica de Segunda Ordem}

A linguagem da lógica de segunda ordem permite quantificar não apenas
sobre !!a{domain} de indivíduos, mas também sobre relações nesse
!!{domain}. Dada uma linguagem de primeira ordem $\Lang{L}$, para cada
$k$ adiciona-se !!{variable}s $R$ que variam sobre relações $k$-árias, e
permite-se quantificação sobre essas variáveis. Se $R$ é !!a{variable} para uma
relação $k$-ária, e $t_1$, \dots,~$t_k$ são termos ordinários (de primeira ordem),
$\Atom{R}{t_1,\dots,t_k}$ é uma !!{formula} atômica. Caso contrário, o conjunto de
!!{formula}s é definido exatamente como no caso da lógica de primeira ordem, com
cláusulas adicionais para quantificação de segunda ordem. Note que temos apenas
a !!{identity} para termos de primeira ordem: se $R$ e $S$ são !!{variable}s de relação
de mesma aridade~$k$, podemos definir $\eq[R][S]$ como uma
abreviação de
\[
\lforall[x_1 \dots][\lforall[x_k][(\Atom{R}{x_1, \dots, x_k} \liff
  \Atom{S}{x_1, \dots, x_k})]].
\]

As regras para a lógica de segunda ordem simplesmente estendem as regras de quantificadores
às novas variáveis de segunda ordem. Aqui, porém, é preciso ter um pouco
de cuidado para explicar como essas variáveis interagem com os !!{predicate}s
de~$\Lang{L}$, e com !!{formula}s de~$\Lang{L}$ de modo mais geral. No mínimo,
variáveis de relação contam como termos, de modo que temos inferências da
forma
\[
!A(R) \Proves \lexists[R][!A(R)]
\]
Mas se $\Lang{L}$ é a linguagem da aritmética com um símbolo de relação
constante~$<$, também esperaríamos que a seguinte inferência fosse
válida:
\[
x < y \Proves \lexists[R][\Atom{R}{x,y}]
\]
ou, para uma dada !!{formula}~$!A$,
\[
!A(x_1, \dots, x_k) \Proves \lexists[R][\Atom{R}{x_1,\dots,x_k}]
\]
Mais geralmente, podemos querer permitir inferências da forma
\[
\Subst{!A}{\lambd[\vec x][!B(\vec x)]}{R} \Proves
\lexists[R][!A]
\]
onde $\Subst{!A}{\lambd[\vec x][!B(\vec x)]}{R}$ denota o resultado
de substituir cada !!{formula} atômica da forma
$\Obj{R}{t_1,\dots,t_k}$ em~$!A$ por $!B(t_1, \dots, t_k)$. Esta última
regra é equivalente a ter um {\em esquema de compreensão}, isto é, um
axioma da forma
\[
\lexists[R][\lforall[x_1, \dots, x_k][(!A(x_1, \dots, x_k) \liff
\Atom{R}{x_1, \dots, x_k})]],
\]
um para cada !!{formula}~$!A$ na linguagem de segunda ordem, em que
$R$ não é uma variável livre. (Exercício: mostre que se $R$ for permitida
em~$!A$, esse esquema é inconsistente!)

Quando lógicos se referem aos ``axiomas da lógica de segunda ordem'', geralmente
querem dizer a extensão mínima da lógica de primeira ordem pelas
regras de quantificadores de segunda ordem juntamente com o esquema de
compreensão. Mas é frequentemente interessante estudar subsistemas mais fracos desses
axiomas e regras. Por exemplo, note que em toda a sua generalidade
o esquema de axiomas de compreensão é \emph{impredicativo}: ele permite
afirmar a existência de uma relação $\Atom{R}{x_1, \dots, x_k}$
que é ``definida'' por !!a{formula} com quantificadores de segunda ordem; e
esses quantificadores variam sobre o conjunto de todas essas relações---um conjunto
que inclui $R$ em si!{} Por volta da virada do século XX, uma
reação comum ao paradoxo de Russell foi atribuir a culpa a tais
definições e evitá-las ao desenvolver os fundamentos da
matemática. Se se proíbe o uso de quantificadores de segunda ordem na
!!{formula}~$!A$, tem-se uma forma {\em predicativa} de
compreensão, que é um pouco mais fraca.

Do ponto de vista semântico, pode-se pensar em uma !!{structure} de segunda ordem
como consistindo em uma !!{structure} de primeira ordem para a
linguagem, juntamente com um conjunto de relações no !!{domain} sobre o qual
variem os quantificadores de segunda ordem (mais precisamente, para cada $k$ há
um conjunto de relações de aridade~$k$). É claro que, se a compreensão estiver
incluída no sistema de !!{derivation}, então temos o requisito adicional de que
há relações suficientes na ``parte de segunda ordem'' para satisfazer os
axiomas de compreensão---caso contrário, o sistema de !!{derivation} não é correto!{} Uma
maneira fácil de garantir que há relações suficientes é tomar
a parte de segunda ordem como consistindo de \emph{todas} as relações sobre a
parte de primeira ordem. Uma tal !!{structure} é chamada \emph{plena}, e, de certo
modo, é realmente a ``!!{structure} pretendida'' para a linguagem. Se
restringirmos a atenção a !!{structure}s plenas, temos o que se conhece
como a semântica de segunda ordem \emph{plena}. Nesse caso, especificar uma
!!{structure} reduz-se a especificar a parte de primeira ordem, pois o
conteúdo da parte de segunda ordem segue disso implicitamente.

Para resumir, há certa ambiguidade ao falar de lógica de segunda ordem.
Em termos do sistema de !!{derivation}, pode-se ter em mente
\begin{enumerate}
\item Um sistema de !!{derivation} de segunda ordem ``mínimo'', juntamente com alguns
  axiomas de compreensão.
\item O sistema de !!{derivation} de segunda ordem ``padrão'', com compreensão
  plena.
\end{enumerate}
Em termos da semântica, pode interessar
\begin{enumerate}
\item A semântica ``fraca'', em que !!a{structure} consiste em uma parte de primeira ordem,
  juntamente com uma parte de segunda ordem grande o suficiente para satisfazer os
  axiomas de compreensão.
\item A semântica de segunda ordem ``padrão'', em que se consideram
  apenas !!{structure}s plenas.
\end{enumerate}
Quando lógicos não especificam o sistema de !!{derivation} ou a semântica que
têm em mente, geralmente se referem ao segundo item de cada
lista. A vantagem de usar essa semântica é que, como veremos,
ela nos dá descrições categóricas de muitas !!{structure}s matemáticas naturais; ao mesmo tempo, o sistema de !!{derivation} é bastante forte, e
correto para essa semântica. A desvantagem é que o sistema de !!{derivation} \emph{não} é
completo para a semântica; de fato, \emph{nenhum} sistema de !!{derivation} efetivamente
dado é completo para a semântica plena de segunda ordem. Por
outro lado, veremos que o sistema de !!{derivation} \emph{é} completo
para a semântica enfraquecida; isso implica que, se uma sentença não for
demonstrável, então há \emph{alguma} !!{structure}, não necessariamente a
plena, na qual ela é falsa.

A linguagem da lógica de segunda ordem é bastante rica. Pode-se identificar
relações unárias com subconjuntos do !!{domain}, e assim, em particular, pode-se
quantificar sobre esses conjuntos; por exemplo, pode-se expressar indução
para os números naturais com um único axioma
\[
\lforall[R][((\Atom{R}{\Obj{0}} \land \lforall[x][(\Atom{R}{x} \lif
    \Atom{R}{x'})]) \lif \lforall[x][\Atom{R}{x}])].
\]
Se se toma a linguagem da aritmética com os símbolos $\Obj 0, \Obj \prime, +,
\times$ e $<$, pode-se adicionar os seguintes axiomas para descrever seu
comportamento:
\begin{enumerate}
\item $\lforall[x][\lnot x' = \Obj 0]$
\item $\lforall[x][\lforall[y][(s(x) = s(y) \lif x = y)]]$
\item $\lforall[x][(x + \Obj 0) = x]$
\item $\lforall[x][\lforall[y][(x + y') = (x + y)']]$
\item $\lforall[x][(x \times \Obj 0) = \Obj 0]$
\item $\lforall[x][\lforall[y][(x \times y') = ((x \times y) + x)]]$
\item $\lforall[x][\lforall[y][(x < y \liff \lexists[z][y = (x + z')])]]$
\end{enumerate}
Não é difícil mostrar que esses axiomas, juntamente com o axioma
de indução acima, fornecem uma descrição categórica da
!!{structure}~$\Struct{N}$, o modelo padrão da aritmética, desde que
estejamos usando a semântica plena de segunda ordem. Dada qualquer
!!{structure}~$\Struct{M}$ na qual esses axiomas são verdadeiros, defina uma
função~$f$ de $\Nat$ para o !!{domain} de $\Struct{M}$ usando
recursão ordinária em $\Nat$, de modo que $f(0) = \Assign{\Obj 0}{M}$ e
$f(x+1) = \Assign{\prime}{M}(f(x))$. Usando indução ordinária
em~$\Nat$ e o fato de que os axiomas (1) e~(2) valem em~$\Struct M$, vemos
que $f$ é !!{injective}. Para ver que $f$ é !!{surjective},
seja~$P$ o conjunto de elementos de~$\Domain{M}$ que estão no alcance
de~$f$. Como $\Struct M$ é plena, $P$ está no !!{domain} de segunda ordem.
Pela construção de~$f$, sabemos que $\Assign{\Obj
  0}{M}$ está em~$P$, e que $P$~é fechado sob
$\Assign{\prime}{M}$. O fato de que o axioma de indução vale em
$\Struct M$ (em particular, para~$P$) garante que $P$~é igual a
todo o !!{domain} de primeira ordem de~$\Struct M$. Isso mostra que $f$
é !!a{bijection}. Mostrar que $f$ é um homomorfismo não é mais
difícil, usando indução ordinária em $\Nat$ repetidamente.

Em termos teórico-conjuntistas, uma função é apenas um tipo especial de relação;
por exemplo, uma função unária~$f$ pode ser identificada com uma relação
binária~$R$ que satisfaz $\lforall[x][\lexists![y][R(x,y)]]$. Como resultado, pode-se
também quantificar sobre funções. Usando a semântica plena, pode-se
então definir a classe de !!{structure}s infinitas como a classe de
!!{structure}s $\Struct M$ para as quais há uma função injetiva do
!!{domain} de $\Struct M$ para um subconjunto próprio de si mesma:
\[
\lexists[f][(\lforall[x][\lforall[y][(\eq[f(x)][f(y)] \lif \eq[x][y])]]
  \land \lexists[y][\lforall[x][\eq/[f(x)][y]]])].
\]
A negação dessa sentença define então a classe de !!{structure}s
finitas.

Além disso, pode-se definir a classe de ordens bem-fundadas, adicionando o
seguinte à definição de uma ordenação linear:
\[
\lforall[P][(\lexists[x][\Atom{P}{x}] \lif \lexists[x][(\Atom{P}{x}
    \land \lforall[y][(y < x \lif \lnot \Atom{P}{y})])])].
\]
Isso afirma que todo conjunto não vazio tem um elemento mínimo, módulo a
identificação de ``conjunto'' com ``relação unária''. Como outro
exemplo, pode-se expressar a noção de conexidade para grafos, dizendo
que não há separação não trivial dos vértices em partes desconectadas:
\[
\lnot \lexists[A][(\lexists[x][A(x)] \land \lexists[y][\lnot A(y)]
  \land \lforall[w][\lforall[z][((\Atom{A}{w} \land \lnot \Atom{A}{z})
      \lif \lnot \Atom{R}{w,z})]])].
\]
Como mais um exemplo, pode-se tentar, como exercício, definir a
classe de !!{structure}s finitas cujo !!{domain} tem tamanho par. De
modo ainda mais impressionante, pode-se fornecer uma descrição categórica dos números
reais como um corpo ordenado completo que contém os racionais.

Em suma, a lógica de segunda ordem é muito mais expressiva do que a lógica de primeira ordem.
Essa é a boa notícia; agora a má. Já
mencionamos que não há sistema de !!{derivation} efetivo que seja completo para
a semântica plena de segunda ordem. Para o bem ou para o mal, muitas das
propriedades da lógica de primeira ordem estão ausentes, incluindo compacidade e
os teoremas de L\"owenheim--Skolem.

Por outro lado, se se está disposto a abrir mão da semântica plena de segunda ordem
em favor da mais fraca, então o sistema de !!{derivation} mínimo de segunda ordem
é completo para essa semântica. Em outras palavras, se
lemos $\Proves$ como ``demonstra no sistema mínimo'' e $\Entails$ como
``implica logicamente na semântica fraca'', podemos mostrar que
sempre que $\Gamma \Entails !A$, então $\Gamma \Proves !A$. Se se quer
incluir axiomas específicos de compreensão no sistema de !!{derivation}, é preciso
restringir a semântica a !!{structure}s de segunda ordem que satisfaçam
esses axiomas: por exemplo, se $\Delta$ consiste em um conjunto de
axiomas de compreensão (possivelmente todos eles), temos que, se $\Gamma
\cup \Delta \Entails !A$, então $\Gamma \cup \Delta \Proves !A$. Em
particular, se $!A$ não for demonstrável usando os axiomas de compreensão que
estamos considerando, então há um modelo de $\lnot !A$ no qual esses
axiomas de compreensão ainda assim valem.

A maneira mais fácil de ver que o teorema de completude vale para a
semântica fraca é pensar na lógica de segunda ordem como uma lógica de múltiplas ordenações,
da seguinte forma. Uma ordenação é interpretada como o !!{domain} ordinário
de ``primeira ordem'', e então, para cada $k$, temos !!a{domain}
de ``relações de aridade $k$.'' Tomamos a linguagem com símbolos de relação
incorporados ``$\Atom{\Obj{true}_k}{R,x_1,\dots,x_k}$'', destinados a
afirmar que $R$ vale de $x_1$, \dots,~$x_k$, onde $R$ é uma
variável da ordenação ``relação $k$-ária'' e $x_1$, \dots,~$x_k$ são
objetos da ordenação de primeira ordem.

Com essa identificação, a semântica fraca de segunda ordem é essencialmente
a semântica usual da lógica de múltiplas ordenações; e já
observamos que a lógica de múltiplas ordenações pode ser embutida na lógica de primeira ordem.
Módulo as traduções de ida e volta, então, a concepção mais fraca de
lógica de segunda ordem é realmente uma forma disfarçada de lógica de primeira ordem,
em que o !!{domain} contém tanto ``objetos'' quanto ``relações''
governados pelos axiomas apropriados.

\end{document}
